\section{Waymo：系统工程与工程师思维}

博士毕业后，他选择自动驾驶，因为它是当时最清晰的 physical world AI 形态。他在 Waymo 的学习重点，不只是自动驾驶算法，而是自动驾驶整套系统如何工作：感知、定位、离线高精地图、预测、决策、规划、控制、仿真和云端工具链如何组合成一个可靠系统。

\subsection{自动驾驶架构的历史连续性}

他回看 2008 年前后的自动驾驶论文后发现，2018 年左右的总体架构和早期框架并无根本差异：系统仍然被拆成感知、定位、地图、规控等模块。真正变化的是 AI 技术成熟后，许多感知模块从 clustering、规则和传统方法逐步替换成神经网络。

这带来一个重要区分：传统自动驾驶大架构的底层逻辑更接近 robotics，强调模块化、可解释和 corner case；AI native 路线更强调数据驱动、端到端和整体 benchmark 改善。两者不是简单谁对谁错，而是在可靠性、可迭代性、可解释性和规模化之间做不同取舍。

\begin{warningbox}{模块化不是落后，端到端也不是万能}
模块化系统更容易定位 corner case 和责任边界，但可能转向慢、局部规则堆积；端到端系统更容易用数据提升整体指标，但可能引入不可解释失败。物理世界 AI 常常需要在两者之间重新设计边界。
\end{warningbox}

\subsection{工程师思维：拆解与测量}

高继扬总结 Waymo 给他的核心训练是工程师思维：拆解加测量。把复杂问题拆成可处理的子问题，再拆到代码和测试；同时从顶层指标、中间指标到底层单元测试建立测量体系。这和物理竞赛的解题不同，后者更多是题型映射，工程系统则要持续运转、持续定位问题、持续回到总体目标。

\begin{importantbox}{工程师思维}
工程师思维不是“会写代码”，而是把复杂系统拆成可验证模块，并用指标体系把局部改动和顶层结果连接起来。
\end{importantbox}

\subsection{本章小结}

Waymo 给高继扬的是完整系统观和工程训练：知道物理世界 AI 系统如何拆，如何测，如何长期执行。但它离客户、产品和公司经营较远，因此他下一步选择去更接近量产交付的 Momenta。

